\documentclass[a4paper,12pt]{article}
\usepackage[T2A]{fontenc}
\usepackage[utf8]{inputenc}
\usepackage[russian]{babel}
\usepackage{amsmath,amssymb,mathtools}
\usepackage{icomma}
\usepackage{geometry}
\geometry{left=22mm,right=22mm,top=21mm,bottom=22mm}
\usepackage{microtype}
\renewcommand{\familydefault}{\sfdefault}
\usepackage{tikz}
\usetikzlibrary{arrows.meta,positioning,shapes.geometric}
\usepackage{pgfplots}
\pgfplotsset{compat=1.18}
\usepackage{tabularx,booktabs}
\usepackage{enumitem}
\usepackage{fancyhdr}
\usepackage{parskip}
\usepackage{xcolor}
\usepackage{hyperref}
\definecolor{accent}{HTML}{008080}
\definecolor{darkaccent}{HTML}{005456}
\definecolor{textgray}{HTML}{3E4546}
\definecolor{softgray}{HTML}{D6E3E2}
\hypersetup{colorlinks=true,linkcolor=darkaccent,urlcolor=darkaccent,
 pdftitle={Чек-лист: обратимые и взаимно обратные функции},
 pdfauthor={Лёвин Артём Александрович}}
\linespread{1.17}
\setlength{\parindent}{0pt}
\setlength{\parskip}{5pt}
\setlist[itemize]{leftmargin=6mm,topsep=3pt,itemsep=3pt}
\setlist[enumerate]{leftmargin=8mm,topsep=4pt,itemsep=8pt}
\renewcommand{\arraystretch}{1.23}
\pagestyle{fancy}
\fancyhf{}
\fancyhead[L]{\footnotesize\color{textgray}Обратимые функции}
\fancyhead[R]{\footnotesize\color{textgray}Анна Трапезникова}
\fancyfoot[C]{\footnotesize\color{textgray}\thepage}
\renewcommand{\headrulewidth}{0.25pt}
\setlength{\headheight}{15pt}
\newcommand{\sect}[1]{\vspace{5pt}{\large\bfseries\color{darkaccent}#1}\par\vspace{3pt}}
\newcommand{\idea}[1]{\textcolor{darkaccent}{\textbf{Ключевая идея.}} #1}
\newcommand{\warn}[1]{\textcolor{darkaccent}{\textbf{Важно.}} #1}
\newcommand{\example}[1]{\textbf{Пример.} #1}
\newcommand{\chk}{\(\square\)\,}
\newcommand{\R}{\mathbb{R}}
\begin{document}
\thispagestyle{empty}
\begin{center}
\vspace*{26mm}
{\large\color{darkaccent}МАТЕМАТИКА}\par
\vspace{19mm}
{\LARGE\color{textgray}Чек-лист}\par
\vspace{10mm}
{\LARGE\color{accent}Обратимые и взаимно обратные функции}\par
\vspace{5mm}
{\large\color{accent}Область определения, графики, композиции}\par
\vspace{17mm}
{\normalsize\color{textgray}Алгебра · 10 класс · углублённый уровень}\par
\vspace{34mm}
{\large Анна Трапезникова}\par
\vspace{13mm}
{\normalsize 15 ноября 2026 г.}\par
\vfill
{\small Лёвин Артём Александрович}\par
{\small эксклюзивно для Анны Трапезниковой}\par
\vspace{19mm}
\begin{tikzpicture}[x=1cm,y=1cm]
\draw[accent,line width=.9pt] (0,0) .. controls (1.4,.56) and (2.7,-.53) .. (4.1,0)
 .. controls (5.5,.55) and (6.8,-.55) .. (8.2,0);
\fill[accent] (4.1,0) circle (1.35pt);
\end{tikzpicture}
\end{center}
\clearpage

\sect{1. Что нужно уметь}
\begin{itemize}
\item \chk Определять, является ли функция \textbf{обратимой} на заданной области.
\item \chk Находить обратную функцию: менять роли \(x\) и \(y\), выражать переменную.
\item \chk Правильно записывать \(D(f^{-1})\) и \(E(f^{-1})\).
\item \chk Строить и узнавать графики взаимно обратных функций.
\item \chk Учитывать ограничения при квадратных корнях и чётных степенях.
\item \chk Проверять взаимную обратность через \(f(g(x))\) и \(g(f(x))\).
\end{itemize}

\sect{2. Что означает «обратимая функция»}
Функция \(y=f(x)\) \textbf{обратима на своей области определения}, если каждое принимаемое ею значение \(y\) достигается \emph{ровно при одном} допустимом значении \(x\).

\idea{Достаточно мысленно провести горизонтальную прямую: для обратимой функции она пересечёт график не более одного раза.}

\begin{center}
\begin{tabularx}{\linewidth}{@{}p{37mm}X@{}}
\toprule
Функция и область & Обратима? Почему?\\
\midrule
\(f(x)=3x+5\), \(x\in\R\) & Да. Разным \(x\) соответствуют разные значения \(y\).\\
\(f(x)=x^3\), \(x\in\R\) & Да. Кубическая функция строго возрастает.\\
\(f(x)=x^2\), \(x\in\R\) & Нет. Например, \(f(-2)=f(2)=4\).\\
\(f(x)=x^2\), \(x\ge0\) & Да. На правой ветви параболы каждое \(y\) достигается один раз.\\
\bottomrule
\end{tabularx}
\end{center}

\warn{У любой функции одному \(x\) соответствует одно \(y\). Для \emph{обратимости} требуется обратное свойство: одному значению \(y\) соответствует единственное \(x\).}

\sect{3. Как найти обратную функцию}
Пусть \(f(x)=3x+5\). Запишем \(y=3x+5\) и выразим \(x\):
\[
3x=y-5\quad\Longrightarrow\quad x=\frac{y-5}{3}.
\]
Теперь поменяем обозначения входа и выхода:
\[
\boxed{f^{-1}(x)=\frac{x-5}{3}.}
\]
\textbf{Правило записи:} сначала укажите исходную функцию \(f(x)\), затем её обратную \(f^{-1}(x)\) или \(g(x)\). Символ \(f^{-1}\) \emph{не означает} \(1/f\).

\example{Линейная функция с отрицательным коэффициентом:}
\[
y=-2x+1\ \Longrightarrow\ 2x=1-y\ \Longrightarrow\
x=\frac{1-y}{2}\ \Longrightarrow\
\boxed{f^{-1}(x)=\frac{1-x}{2}.}
\]
Для любой линейной функции \(f(x)=ax+b\), \(a\ne0\), обратная существует на всей \(\R\):
\[
\boxed{f^{-1}(x)=\frac{x-b}{a}.}
\]

\clearpage
\sect{4. Графики: отражение относительно \(y=x\)}
При переходе к обратной функции координаты каждой точки меняются местами:
\[
(a;b)\ \longleftrightarrow\ (b;a).
\]
Поэтому графики \(y=f(x)\) и \(y=f^{-1}(x)\) \textbf{симметричны относительно прямой \(y=x\)}. Прямая \(y=x\) проходит через I и III координатные четверти.

\begin{center}
\begin{tikzpicture}
\begin{axis}[
 width=12.5cm,height=9cm,axis lines=middle,
 xmin=-4.1,xmax=4.1,ymin=-4.1,ymax=4.1,
 xminorticks=false,yminorticks=false,
 xtick={-4,-3,-2,-1,0,1,2,3,4},
 ytick={-4,-3,-2,-1,0,1,2,3,4},
 xlabel={$x$},ylabel={$y$},
 xlabel style={at={(axis description cs:1,0.48)},anchor=west},
 ylabel style={at={(axis description cs:0.51,1)},anchor=south},
 axis equal image,grid=major,grid style={softgray!55},
 legend style={at={(0.02,0.98)},anchor=north west,draw=none,fill=white,font=\small},
 ticklabel style={font=\scriptsize},clip=false]
\addplot[accent,very thick,domain=-1.55:2.5,samples=2] {2*x-1};
\addlegendentry{$f(x)=2x-1$}
\addplot[darkaccent,very thick,densely dotted,domain=-4:4,samples=2] {(x+1)/2};
\addlegendentry{$f^{-1}(x)=\frac{x+1}{2}$}
\addplot[gray,dashed,domain=-4:4,samples=2] {x};
\addlegendentry{$y=x$}
\addplot[only marks,mark=*,mark size=2.3pt] coordinates {(0,-1) (-1,0) (1,1)};
\end{axis}
\end{tikzpicture}
\end{center}
Например, точке \((0;-1)\) графика \(f\) соответствует точка \((-1;0)\) графика \(f^{-1}\).

\warn{Симметрия относительно \(y=x\) и симметрия относительно начала координат --- разные преобразования. Для взаимно обратных функций используется именно отражение относительно \(y=x\).}

\sect{5. Область определения и множество значений}
Обратная функция меняет местами множество допустимых аргументов и множество получаемых значений:
\[
\boxed{D(f^{-1})=E(f),\qquad E(f^{-1})=D(f).}
\]
У функции \(f(x)=2x-1\) оба множества равны \(\R\); у обратной \(f^{-1}(x)=(x+1)/2\) --- тоже.

\clearpage
\sect{6. Почему параболе нужно ограничение}
На всей числовой прямой функция \(f(x)=x^2\) необратима: числу \(9\) соответствуют \(x=-3\) и \(x=3\). Если оставить одну ветвь, появляется обратная функция:
\begin{center}
\small
\begin{tabularx}{\linewidth}{@{}p{31mm}X X@{}}
\toprule
Исходная область & Исходная функция & Обратная функция\\
\midrule
\(x\ge0\) & \(f(x)=x^2\) & \(f^{-1}(x)=\sqrt{x}\), \(x\ge0\)\\
\(x\le0\) & \(f(x)=x^2\) & \(f^{-1}(x)=-\sqrt{x}\), \(x\ge0\)\\
\bottomrule
\end{tabularx}
\end{center}
В обоих случаях множество значений исходной функции равно \([0;+\infty)\), поэтому область определения обратной тоже \([0;+\infty)\). Множество значений обратной зависит от выбранной ветви: соответственно \([0;+\infty)\) или \((-\infty;0]\).

\sect{7. Квадратный корень и потерянный знак}
\example{Найдём функцию, обратную к \(f(x)=-\sqrt{x}\), \(x\ge0\).}
Из \(y=-\sqrt{x}\) следует \(y\le0\). Возведём в квадрат:
\[
y^2=x\quad\Longrightarrow\quad
\boxed{f^{-1}(x)=x^2,\quad x\le0.}
\]
Здесь \(D(f)=[0;+\infty)\), \(E(f)=(-\infty;0]\), поэтому
\[
D(f^{-1})=(-\infty;0],\qquad E(f^{-1})=[0;+\infty).
\]
\warn{После возведения в квадрат нельзя забывать условие \(y\le0\). Без него полученная формула \(x^2\) будет иметь лишнюю половину области определения.}

\sect{8. Кубические корни и нечётные степени}
Для всех действительных \(x\) определены \(\sqrt[3]{x}\) и \(\sqrt[5]{x}\), в том числе при отрицательных \(x\). Поэтому:
\[
\begin{aligned}
f(x)&=x^3+1
&&\Longrightarrow& f^{-1}(x)&=\sqrt[3]{x-1},\\
f(x)&=-x^5
&&\Longrightarrow& f^{-1}(x)&=-\sqrt[5]{x}.
\end{aligned}
\]
Дробный показатель с нечётным знаменателем удобно трактовать через вещественный корень:
\[
(\sqrt[5]{x})^3=x^{3/5},\qquad
f^{-1}(x)=(\sqrt[3]{x})^5=x^{5/3}.
\]
Показатели \(3/5\) и \(5/3\) здесь можно использовать и для отрицательных чисел благодаря корням нечётных степеней.

\clearpage
\sect{9. Как проверить взаимную обратность}
Для \(g=f^{-1}\) выполняются два равенства, \emph{каждое на своей области}:
\[
\boxed{f(g(x))=x\quad (x\in D(g)),\qquad
 g(f(x))=x\quad (x\in D(f)).}
\]
\example{Проверим \(f(x)=-x^3\) и \(g(x)=-\sqrt[3]{x}\).}
\[
\begin{aligned}
f(g(x))&=-\bigl(-\sqrt[3]{x}\bigr)^3=x,\\
g(f(x))&=-\sqrt[3]{-x^3}=x.
\end{aligned}
\]
Значит, функции взаимно обратны на \(\R\).

\example{А что, если забыть минус: \(f(x)=-x^5\), \(g(x)=\sqrt[5]{x}\)?}
\[
\begin{aligned}
f(g(x))&=-\bigl(\sqrt[5]{x}\bigr)^5=-x,\\
g(f(x))&=\sqrt[5]{-x^5}=-x.
\end{aligned}
\]
При \(x\ne0\) получаем \(-x\ne x\). Следовательно, функции \textbf{не взаимно обратны}. Правильный вариант --- \(g(x)=-\sqrt[5]{x}\).

\sect{10. Контроль знака: \(\sqrt{x^2}=|x|\)}
Для обратной \(g(x)=x^2\), \(x\le0\), из раздела 7:
\[
\begin{aligned}
g(f(x))&=(-\sqrt{x})^2=x, && x\ge0,\\
f(g(x))&=-\sqrt{x^2}=-|x|=x, && x\le0.
\end{aligned}
\]
\warn{Равенство \(\sqrt{x^2}=x\) справедливо только при \(x\ge0\). В общем случае \(\sqrt{x^2}=|x|\).}

\sect{11. Частые ошибки}
\begin{itemize}
\item \textbf{Перепутано направление соответствия.} Обратимость проверяют по числу аргументов \(x\) для одного значения \(y\).
\item \textbf{Забыто переименование.} Из \(x=(y-5)/3\) в ответе пишут \(f^{-1}(x)=(x-5)/3\).
\item \textbf{Пропущена область определения.} Для \(-\sqrt{x}\) обратная функция \(x^2\) задаётся только при \(x\le0\).
\item \textbf{Потерян минус.} Обратная к \(-x^5\) содержит минус перед корнем пятой степени.
\item \textbf{Неверно раскрыт корень.} \(\sqrt{x^2}=|x|\), поэтому знак аргумента нужно учитывать.
\item \textbf{Проверка только одной композиции.} Полезно проверять обе и указывать соответствующие области.
\end{itemize}

\clearpage
\begin{center}{\LARGE\color{darkaccent}Алгоритм нахождения обратной функции}\par\vspace{9mm}\end{center}
\begin{center}
\begin{tikzpicture}[
 >={Latex[length=2.5mm]},
 term/.style={ellipse,draw=accent,thick,align=center,minimum width=5.4cm,minimum height=.83cm,font=\small},
 action/.style={rounded corners=2pt,draw=darkaccent,thick,align=center,text width=10cm,minimum height=.92cm,inner sep=4pt,font=\small},
 branch/.style={rounded corners=2pt,draw=darkaccent,thick,align=center,text width=4.2cm,minimum height=1.1cm,inner sep=4pt,font=\small},
 decision/.style={diamond,draw=darkaccent,thick,aspect=2.4,align=center,text width=3.1cm,inner sep=1.5pt,font=\small},
 arr/.style={->,draw=darkaccent,thick}
]
\node[term] (start) at (0,0) {Дана функция $y=f(x)$ и её область $D(f)$};
\node[decision] (q) at (0,-2.0) {Функция обратима?};
\node[branch] (restrict) at (-3.35,-3.95) {Нет: выберите область, на которой значения не повторяются};
\node[branch] (keep) at (3.35,-3.95) {Да: сохраните исходную область определения};
\node[action] (solve) at (0,-6.0) {Запишите $y=f(x)$ и выразите $x$ через $y$};
\node[action] (swap) at (0,-7.42) {Переименуйте переменные: получите $g(x)=f^{-1}(x)$};
\node[action] (domains) at (0,-8.84) {Области меняются местами: $D(g)=E(f)$, $E(g)=D(f)$};
\node[action] (check) at (0,-10.26) {Проверьте $f(g(x))=x$ и $g(f(x))=x$ на своих областях};
\node[term] (done) at (0,-11.72) {Ответ: обратная функция и её область};
\draw[arr] (start) -- (q);
\draw[arr] (q.south west) -| node[pos=.18,above left,font=\scriptsize]{нет} (restrict.north);
\draw[arr] (q.south east) -| node[pos=.18,above right,font=\scriptsize]{да} (keep.north);
\draw[draw=darkaccent,thick] (restrict.south) -- (-3.35,-5.25) -- (0,-5.25);
\draw[draw=darkaccent,thick] (keep.south) -- (3.35,-5.25) -- (0,-5.25);
\draw[arr] (0,-5.25) -- (solve.north);
\draw[arr] (solve) -- (swap);
\draw[arr] (swap) -- (domains);
\draw[arr] (domains) -- (check);
\draw[arr] (check) -- (done);
\end{tikzpicture}
\end{center}

\clearpage
\begin{center}{\LARGE\color{darkaccent}Тренировочный блок}\par\vspace{3mm}
{\normalsize Путь А · 10 заданий · от простого к сложному}\end{center}
\vspace{4mm}
В каждом задании указывайте область определения обратной функции. Если требуется проверить взаимную обратность, учитывайте области обеих функций.

\textbf{Базовый уровень}
\begin{enumerate}
\item Найдите обратную к функции \(f(x)=4x-7\), \(x\in\R\).
\item Найдите обратную к функции \(f(x)=-5x+4\), \(x\in\R\).
\item Найдите обратную к функции \(f(x)=\dfrac{3x-1}{2}\), \(x\in\R\).
\item Найдите обратную к функции \(f(x)=x^3+8\), \(x\in\R\).
\end{enumerate}

\textbf{Средний уровень}
\begin{enumerate}[resume]
\item Обратима ли \(f(x)=x^2-4\) на \(\R\)? Ограничьте её область условием \(x\ge0\), найдите обратную функцию, её \(D\) и \(E\).
\item Для \(f(x)=1-\sqrt{x+2}\) найдите исходные \(D(f)\) и \(E(f)\), обратную функцию и \(D(f^{-1})\), \(E(f^{-1})\).
\item Найдите обратную к \(f(x)=\dfrac{x+1}{2x-3}\), \(x\ne\frac32\), и укажите её область определения и множество значений.
\end{enumerate}

\textbf{Повышенный уровень}
\begin{enumerate}[resume]
\item Являются ли взаимно обратными \(f(x)=-x^5\) и \(g(x)=\sqrt[5]{x}\) на \(\R\)? Проверьте с помощью двух композиций и укажите верную обратную к \(f\).
\item Проверьте, являются ли взаимно обратными \(f(x)=-\sqrt{x}\), \(x\ge0\), и \(g(x)=x^2\), \(x\le0\). Вычислите обе композиции и укажите области \(D\) и множества \(E\).
\item Для функции \(f(x)=(\sqrt[5]{x})^3\), \(x\in\R\), найдите обратную, запишите её через корень и проверьте обе композиции.
\end{enumerate}

\clearpage
\begin{center}{\LARGE\color{darkaccent}Ответы к тренировочному блоку}\par\vspace{7mm}\end{center}
\normalsize
\begin{tabularx}{\linewidth}{@{}p{9mm}X@{}}
\toprule
\textbf{№} & \textbf{Ответ}\\
\midrule
1 & \(f^{-1}(x)=\dfrac{x+7}{4}\), \(D(f^{-1})=\R\).\\
2 & \(f^{-1}(x)=\dfrac{4-x}{5}\), \(D(f^{-1})=\R\).\\
3 & \(f^{-1}(x)=\dfrac{2x+1}{3}\), \(D(f^{-1})=\R\).\\
4 & \(f^{-1}(x)=\sqrt[3]{x-8}\), \(D(f^{-1})=\R\).\\
5 & На \(\R\) необратима. При \(x\ge0\): \(f^{-1}(x)=\sqrt{x+4}\), \(D(f^{-1})=[-4;+\infty)\), \(E(f^{-1})=[0;+\infty)\).\\
6 & \(D(f)=[-2;+\infty)\), \(E(f)=(-\infty;1]\); \(f^{-1}(x)=(1-x)^2-2\), \(D(f^{-1})=(-\infty;1]\), \(E(f^{-1})=[-2;+\infty)\).\\
7 & \(f^{-1}(x)=\dfrac{3x+1}{2x-1}\), \(D(f^{-1})=\R\setminus\{\frac12\}\), \(E(f^{-1})=\R\setminus\{\frac32\}\).\\
8 & Нет: \(f(g(x))=g(f(x))=-x\); верная обратная \(f^{-1}(x)=-\sqrt[5]{x}\).\\
9 & Да: \(g(f(x))=x\) при \(x\ge0\), \(f(g(x))=x\) при \(x\le0\). Для \(f\): \(D=[0;+\infty)\), \(E=(-\infty;0]\); для \(g\) --- наоборот.\\
10 & \(f^{-1}(x)=(\sqrt[3]{x})^5\), \(D(f^{-1})=\R\); обе композиции равны \(x\) на \(\R\).\\
\bottomrule
\end{tabularx}
\normalsize
\vspace{12mm}
\sect{Краткая памятка}
\begin{tabularx}{\linewidth}{@{}p{0.40\linewidth}X@{}}
\toprule
Сигнал в задаче & Что сделать\\
\midrule
Вопрос об обратимости & Убедитесь: каждое значение \(y\) соответствует одному \(x\).\\
Нужно найти обратную & Выразите \(x\) через \(y\), затем поменяйте имена переменных.\\
Функция \(x^2\) & Укажите выбранную ветвь: \(x\ge0\) или \(x\le0\).\\
Корень чётной степени & Проверьте знак, ОДЗ и \(\sqrt{x^2}=|x|\).\\
График обратной & Отразите относительно \(y=x\).\\
Проверка пары & Рассчитайте \(f(g(x))\) и \(g(f(x))\) с учётом областей.\\
\bottomrule
\end{tabularx}
\end{document}
